% ==============================================================================
% ================================ Zusammenfassung ==============================
% ==============================================================================

\chapter*{Zusammenfassung}
\addcontentsline{toc}{chapter}{Zusammenfassung}
\label{chap:Zusammenfassung}
\thispagestyle{fancy}
\section*{Wissenschaftliche Reflexion (\textit{Agilität zwischen Struktur und Psychologie})} \par Die Analyse des Rollenwechsels im Projektmanagement verdeutlicht, dass Agilität weit über ein bloßes methodisches Toolset hinausgeht. Sie stellt primär einen tiefgreifenden kulturellen \textbf{Mindset-Shift} dar. Ein zentrales Erkenntnismoment liegt in der Bedeutung des sogenannten „\textbf{\textit{Leadership Gap}}“: Erfolg in agilen Systemen resultiert nicht aus der Maximierung von Kontrolle, sondern aus dem bewussten Mut zur Lücke. Erst wenn der Scrum Master als Servant Leader den Raum für Selbstorganisation freigibt, kann sich eine belastbare \textbf{psychologische Sicherheit} im Team entfalten. \par Wissenschaftliches Arbeiten bedeutet in diesem Kontext, die intuitive Begeisterung für flache Hierarchien durch objektive Kriterien wie \textbf{Validität} und \textbf{Reliabilität} abzusichern. Das Resümee dieser Untersuchung unterstreicht, dass in einer \acs{VUCA}-Welt die Transformation zum Scrum Master die Abkehr von der „\textbf{\textit{Fachkompetenzfalle}}“ hin zu einer moderierenden Führung markiert. Letztlich entscheiden die \textbf{Soft Facts} über die Belastbarkeit der Projektergebnisse, denn Prozesse bieten lediglich den Rahmen, während Menschen den Erfolg gestalten. 

% ==============================================================================
% ================================ Zusammenfassung ==============================
% ==============================================================================